% TEMPLATE-MILC-v3.tex - plantilla maestra UNICA de las guias MILC v3.
%
% La usan las cuatro familias de guias, cada una con su driver:
%   · Grados 6-11          -> scripts/build-guias-g11.py
%   · Bebras               -> scripts/build-guias-bebras.py
%   · Semillero            -> scripts/build-guias-semillero.py
%   · Territorio interior  -> scripts/build-guias-territorio-interior.py
%
% Antes habia tres copias casi identicas (~400 lineas iguales, 6 distintas):
% cada mejora habia que aplicarla tres veces, y en la practica no se hacia
% (el motor de recursos vivia solo en dos de ellas). Lo que varia entre
% familias esta parametrizado en 6 placeholders que cada driver llena
% (se nombran aqui SIN su sintaxis real: el builder reemplaza por texto plano,
% tambien dentro de los comentarios, y un valor multilinea rompe el preambulo):
%
%   HEADER_IZQ        encabezado izquierdo de pagina
%   HEADER_DER        encabezado derecho de pagina
%   PORTADA_ETIQUETA  rotulo sobre el numero grande (GRADO/MOMENTO/MODULO)
%   PORTADA_SERIE     linea de serie de la portada
%   PORTADA_ID        identificador de la guia en la portada
%   PORTADA_CIRCULO   el circulito de grado (°), o vacio si no aplica
%   PDF_TITULO        titulo en texto plano (sin \textbf ni comillas TeX) para
%                     los metadatos del PDF (/Title); lo produce lib_xelatex.texto_pdf
%
% Composicion (auditoria 2026-09-03): las bandas de estacion piden espacio
% (\Needspace*) y prohiben el salto justo despues (\nopagebreak), asi no quedan
% huerfanas al pie; las cajas largas (softbox, drawbox, infoband, puentebox) son
% `breakable`, asi una caja de media pagina ya no arrastra todo a la siguiente
% dejando la anterior al 40 %. Todo texto sobre mostaza o verde va en milcNegro
% (blanco sobre #E5B400 da 1,93:1 y sobre #58A923 2,95:1; ninguno pasa WCAG AA).
% El resto de claves las documenta content/guias/_SCHEMA.md. El script valida
% que no queden placeholders sin reemplazar antes de escribir el archivo final.

% Etiquetado PDF (latex-lab): /Lang, estructura y texto alternativo de las
% figuras, para lectores de pantalla. Debe ir ANTES de \documentclass.
\DocumentMetadata{lang=es-CO,tagging=on}
\documentclass[11pt,letterpaper]{article}
% headheight=24pt: la cabecera derecha lleva dos lineas (identidad de la guia
% y estacion actual). headsep se acorta en la misma medida para que el bloque
% de texto quede exactamente donde estaba con headheight=12pt.
\usepackage[margin=2.0cm,headheight=24pt,headsep=13pt]{geometry}
\usepackage[table]{xcolor}
\usepackage{fontspec}
\usepackage[spanish]{babel}
\usepackage{tikz}
% Librerías TikZ para los diagramas de las guías (bloque `recursos:`):
%  · babel  → babel en español vuelve activos caracteres como «>», y TikZ los usa
%             en sus opciones (>=stealth, ->). Sin esta librería, cualquier
%             diagrama con flechas revienta con «\language@active@arg> has an
%             extra }». Debe ir ANTES de que se dibuje nada.
%  · positioning        → sintaxis «below left=2pt and 3pt».
%  · decorations.pathreplacing → llaves ({brace}) para señalar rangos.
\usetikzlibrary{babel,positioning,decorations.pathreplacing}
\usepackage{graphicx}
% Circuitos: el trabajo de electrónica se hace hoy con micro:bit y sus sensores
% integrados (MakeCode, sin cableado), así que no se necesitan esquemáticos.
% Si algún día entran montajes con protoboard, basta descomentar esta línea:
% los diagramas .tex/.tikz declarados en `recursos:` ya se incrustan solos.
% \usepackage{circuitikz}
\usepackage[most]{tcolorbox}
\usepackage{tabularx}
\usepackage{array}
\usepackage{fancyhdr}
\usepackage{titlesec}
\usepackage[document]{ragged2e}  % bandera a la derecha en todo el cuerpo
\usepackage{enumitem}
\usepackage{microtype}
\usepackage{etoolbox}
\usepackage{needspace}
% hyperref al final del preambulo: metadatos del PDF (titulo, autor, idioma,
% familia), marcadores por estacion (\pdfbookmark) y enlaces sin recuadro.
\usepackage[hidelinks,
  pdftitle={Preguntar antes de calcular — phronesis con datos},
  pdfauthor={PhD. Álvaro Cárdenas Orozco},
  pdfsubject={Tecnología e Informática · Grado 8 · Guía 1 · MILC v3},
  pdflang={es-CO},
  bookmarksopen=true,bookmarksnumbered=false]{hyperref}

% Helvetica Neue no trae la flecha U+2192, asi que cada «→» escrito literalmente
% en el YAML se compilaba a NADA, en silencio (462 flechas en 61 guias: rutas del
% tipo «paleta → tipografia → lockup» salian rotas). Mapeamos el caracter a su
% version matematica, que si tiene glifo. Asi el autor puede seguir escribiendo
% «→» comodamente en el YAML.
\usepackage{newunicodechar}
\newunicodechar{→}{\ensuremath{\rightarrow}}
% Mismo problema con los simbolos de marca y vinetas: el «✓» de «una palomita ✓»
% salia como cuadrito vacio en el PDF de todas las guias que lo usaban.
\usepackage{pifont}
\newunicodechar{✓}{\ding{51}}
\newunicodechar{✗}{\ding{55}}
\newunicodechar{✔}{\ding{52}}
\newunicodechar{•}{\textbullet}
\newunicodechar{≥}{\ensuremath{\geq}}
\newunicodechar{≤}{\ensuremath{\leq}}

\setmainfont{Helvetica Neue}
\setsansfont{Helvetica Neue}
\setlength{\parindent}{0pt}
\setlength{\parskip}{6pt}
% Sin particion de palabras: en 6.º-7.º una palabra cortada por guion al final
% de linea («Comuni-cacion») frena la lectura mucho mas que una linea corta.
\hyphenpenalty=10000
\exhyphenpenalty=10000
\pretolerance=2000
\tolerance=3000
\setlength{\emergencystretch}{3em}
\linespread{1.10}
\renewcommand{\arraystretch}{1.25}
\setlist{leftmargin=5mm,itemsep=1mm,topsep=1mm}
\newcolumntype{Y}{>{\RaggedRight\arraybackslash}X}

\definecolor{milcMagenta}{HTML}{E6007E}
\definecolor{milcVino}{HTML}{5A0038}
\definecolor{milcTurquesa}{HTML}{0093A5}
\definecolor{milcVerde}{HTML}{58A923}
\definecolor{milcMostaza}{HTML}{E5B400}
\definecolor{milcOcre}{HTML}{B86600}
\definecolor{milcNegro}{HTML}{241816}
\definecolor{milcGris}{HTML}{F7F7F4}
\definecolor{milcLinea}{HTML}{E8E2DD}
\definecolor{milcGradePrimary}{HTML}{FF6600}
\definecolor{milcGradeDark}{HTML}{9A3E00}
\definecolor{milcGradeSoft}{HTML}{FFE4D0}
\definecolor{milcGradeAccent}{HTML}{CFE7FF}
\definecolor{milcGradeText}{HTML}{FFFFFF}
\color{milcNegro}

\pagestyle{fancy}
\fancyhf{}
% Cabecera de dos lineas: identidad de la guia arriba y, a la derecha y en
% gris, la estacion en curso (\markright desde \milcestacion). Antes de la
% primera estacion la segunda linea va vacia; el \strut mantiene la altura
% para que la linea de arriba no baile entre paginas.
\lhead{\small Tecnología e Informática · Grado 8\\[-1pt]{\footnotesize\strut}}
\rhead{\small Guía 1 · MILC v3\\[-1pt]{\footnotesize\strut\color{milcNegro!65}\rightmark}}
\cfoot{\small\thepage}
\renewcommand{\headrulewidth}{0pt}

% Sobre mostaza (#E5B400) y verde (#58A923) el texto va en milcNegro; sobre el
% resto de la paleta, en blanco. Se usa dentro de fonttitle/fontupper, donde un
% \color posterior manda sobre coltitle.
\newcommand{\milctintasobre}[1]{%
  \ifboolexpr{test{\ifstrequal{#1}{milcMostaza}} or test{\ifstrequal{#1}{milcVerde}}}%
    {\color{milcNegro}}{\color{white}}}

\newtcolorbox{titlebox}[1]{
  enhanced,colback=#1,colframe=#1,arc=7mm,boxrule=0pt,
  left=7mm,right=7mm,top=3mm,bottom=3mm,
  before skip=8pt,after skip=8pt,
  fontupper=\sffamily\bfseries\Large\color{white}
}

\newtcolorbox{softbox}[3]{
  enhanced,breakable,pad at break=3mm,
  colback=#2,colframe=#1,borderline west={4pt}{0pt}{#1},
  arc=2mm,boxrule=.4pt,left=4mm,right=4mm,top=3mm,bottom=3mm,
  before skip=6pt,after skip=8pt,title={#3},coltitle=white,
  colbacktitle=#1,fonttitle=\sffamily\bfseries\milctintasobre{#1},fontupper=\RaggedRight,
  attach boxed title to top left={xshift=3mm,yshift=-2mm},
  boxed title style={arc=2mm, boxrule=0pt}
}

\newtcolorbox{drawbox}[2]{
  enhanced,breakable,pad at break=4mm,
  colback=white,colframe=#1,arc=4mm,boxrule=1pt,
  left=5mm,right=5mm,top=4mm,bottom=4mm,before skip=8pt,after skip=8pt,
  title={#2},coltitle=white,colbacktitle=#1,fonttitle=\sffamily\bfseries\milctintasobre{#1},
  fontupper=\RaggedRight,
  attach boxed title to top left={xshift=4mm,yshift=-2mm},
  boxed title style={arc=3mm, boxrule=0pt}
}

\newtcolorbox{infoband}[1]{
  enhanced,breakable,pad at break=2mm,colback=#1!8,colframe=#1!50,arc=2mm,boxrule=.5pt,
  left=4mm,right=4mm,top=2mm,bottom=2mm,before skip=4pt,after skip=4pt,
  fontupper=\small\RaggedRight
}

% Puente narrativo entre secciones (contrato editorial Conéctate).
% Da continuidad: "Acabas de X. Ahora vas a Y porque Z."
% Visualmente: banda izquierda en color del grado, texto en itálica pequeña.
\newtcolorbox{puentebox}{
  enhanced,breakable,pad at break=2mm,colback=milcGradeSoft,colframe=milcGradeDark,
  borderline west={3pt}{0pt}{milcGradeDark},
  arc=0mm,boxrule=0pt,left=5mm,right=4mm,top=2.5mm,bottom=2.5mm,
  before skip=4pt,after skip=8pt,
  fontupper=\itshape\small\color{milcNegro}\RaggedRight
}
% La flecha va en modo matematico a proposito: Helvetica Neue NO tiene el glifo
% U+2192, asi que \textrightarrow se compilaba a NADA («Missing character: There
% is no → in font Helvetica Neue») y el puente salia sin flecha. En modo
% matematico la flecha viene de la fuente matematica y si se dibuja.
\newcommand{\puente}[1]{%
  \begin{puentebox}%
    \textcolor{milcGradeDark}{$\rightarrow$}\hspace{2mm}#1%
  \end{puentebox}%
}

\newcommand{\linead}[1]{\noindent\color{milcLinea}\rule{#1}{0.5pt}\color{milcNegro}}
\newcommand{\respuesta}[1]{\par\vspace{2mm}\foreach \n in {1,...,#1}{\noindent\color{milcLinea}\rule{\linewidth}{0.5pt}\par\vspace{5mm}}\color{milcNegro}}
\newcommand{\checkbox}{\textcolor{milcGradeDark}{\fbox{\phantom{x}}}\hspace{5pt}}

% ─── Listas aireadas para las guías socioemocionales ────────────────────
% Componer las verificaciones como «\checkbox texto\\» deja el texto que salta
% de línea debajo del cuadrito y apelmaza el bloque. Estos entornos alinean la
% continuación y airean los ítems. Los usa Territorio Interior; cualquier
% familia puede hacerlo.
\newlist{checklista}{itemize}{1}
\setlist[checklista]{
  label=\textcolor{milcGradeDark}{\fbox{\phantom{x}}},
  leftmargin=9mm, labelsep=3.2mm, itemsep=2.4mm,
  topsep=2.5mm, parsep=0pt, partopsep=0pt, before=\RaggedRight
}
\newlist{pasoslista}{enumerate}{1}
\setlist[pasoslista]{
  label=\textcolor{milcGradeDark}{\sffamily\bfseries\arabic*.},
  leftmargin=8mm, labelsep=3mm, itemsep=2.8mm,
  topsep=1mm, parsep=0pt, partopsep=0pt, before=\RaggedRight
}

\newcommand{\phasecard}[4]{
\begin{tcolorbox}[enhanced,colback=#3,colframe=#2,arc=3mm,boxrule=.5pt,height=3.05cm,valign=center,halign=center,left=1.5mm,right=1.5mm,top=2mm,bottom=2mm]
{\sffamily\bfseries\fontsize{22}{24}\selectfont\textcolor{#2}{#1}}\par
{\sffamily\bfseries\small #4}
\end{tcolorbox}
}

% ── Piezas del rediseno 2026 (legibilidad 6.º-11.º) ───────────────────
% Banda de estacion: reemplaza el titlebox macizo. Titulo a la izquierda,
% dato practico (tiempo/modalidad) a la derecha, en una sola linea.
% Nunca huerfana: pide 9 lineas libres (o salta de pagina rellenando con
% \vfill) y prohibe el salto justo despues de la banda. Nueve y no seis:
% la caja partible que sigue necesita ~80pt (titulo adosado, relleno y dos
% lineas) para arrancar en la misma pagina; con menos, tcolorbox la manda
% entera a la siguiente y la banda se queda sola. El penalty 10000 va
% en `after`, ANTES del espacio posterior: un `after skip` (aunque sea 0pt)
% es un \vskip tras la caja, y ese glue es un punto de corte legal que un
% \nopagebreak escrito despues de \end{tcolorbox} ya no cierra. Cada banda
% deja marcador PDF y actualiza la cabecera derecha.
\newcounter{milcestacion}
\newcommand{\milcestacion}[2]{%
\Needspace*{9\baselineskip}%
\stepcounter{milcestacion}%
\pdfbookmark[1]{#2}{milcestacion\themilcestacion}%
\markright{#2}%
\begin{tcolorbox}[enhanced,colback=#1,colframe=#1,arc=2mm,boxrule=0pt,
  left=5mm,right=5mm,top=2.6mm,bottom=2.6mm,before skip=11pt,
  after={\par\nopagebreak\vskip7pt}]
{\sffamily\bfseries\fontsize{15}{18}\selectfont\milctintasobre{#1}#2}
\end{tcolorbox}}

% Tarjeta de paso numerada: sustituye las tablas «Pilar / Aplicacion», que
% eran formato de consulta docente, no de estudiante. El numero va en un
% circulo sobre el borde para que la secuencia se lea de un vistazo.
\newcommand{\milcpaso}[3]{%
\Needspace{4\baselineskip}%
\begin{tcolorbox}[enhanced,colback=white,colframe=milcLinea,arc=1.5mm,boxrule=.9pt,
  left=14mm,right=4mm,top=3mm,bottom=3mm,before skip=4pt,after skip=4pt,
  fontupper=\RaggedRight,
  overlay={\node[anchor=north west,circle,fill=#2,text=white,inner sep=0pt,
    minimum size=7.4mm,font=\sffamily\bfseries\small]
    at ([xshift=3.6mm,yshift=-3.2mm]frame.north west) {#1};}]
#3
\end{tcolorbox}}

% Casilla marcable de tamano util con lapiz (la anterior era \fbox{\phantom{x}},
% de alto variable segun el texto que la seguia).
\newcommand{\milcmarca}{\framebox[3.8mm]{\rule{0pt}{3.4mm}}}

% Fila marcable de lista suelta (checks de la estacion 1).
\newcommand{\milccheckrow}[1]{%
\par\vspace{1.8mm}\noindent
\makebox[6.4mm][l]{\raisebox{-.55mm}{\textcolor{milcNegro}{\milcmarca}}}%
\parbox[t]{\dimexpr\linewidth-6.4mm\relax}{\RaggedRight #1}\par}

% Celda marcable dentro de la rubrica (tabla de autoevaluacion). Misma
% sangria francesa, pero pensada para vivir dentro de una columna X.
\newcommand{\milcopcion}[1]{%
\makebox[5.2mm][l]{\raisebox{-.55mm}{\textcolor{milcNegro}{\milcmarca}}}%
\parbox[t]{\dimexpr\linewidth-5.2mm\relax}{\RaggedRight #1}}

% ── Recursos visuales de la guía ──────────────────────────────────────
% Declarados en el bloque `recursos:` del YAML y emitidos por el builder.
% \guiaFigura   → imagen rasterizada o PDF (foto, carta celeste, montaje).
% \guiaDiagrama → fuente TikZ/circuitikz (.tex) incrustada como vector:
%                 es la vía para circuitos y esquemáticos editables.
% [#1] = texto alternativo (alt) · #2 = ruta relativa al .tex · #3 = pie
% El alt llega al PDF etiquetado (/Alt) y lo lee el lector de pantalla.
\newcommand{\guiaFigura}[3][]{%
\begin{center}
\includegraphics[width=0.88\linewidth,height=0.38\textheight,keepaspectratio,alt={#1}]{#2}\\[1.5mm]
{\footnotesize\itshape\textcolor{milcNegro!75}{#3}}
\end{center}
}

\newcommand{\guiaDiagrama}[2]{%
\begin{center}
\input{#1}\\[1.5mm]
{\footnotesize\itshape\textcolor{milcNegro!75}{#2}}
\end{center}
}

\begin{document}
\thispagestyle{empty}

% ── Portada ───────────────────────────────────────────────────────────
% Antes: un rectangulo de color a pagina completa. Se veia bien en pantalla,
% pero en la fotocopiadora del colegio se comia el toner y salia gris sucio.
% Ahora la banda ocupa el tercio superior y el resto va en blanco.
\begin{tikzpicture}[remember picture,overlay]
  \path[fill=milcGradePrimary]
    (current page.north west) rectangle ([yshift=-10.6cm]current page.north east);

  \node[anchor=north west,text=milcGradeText] at ([xshift=2.0cm,yshift=-1.55cm]current page.north west)
    {\sffamily\bfseries\fontsize{12}{14}\selectfont GRADO};
  \node[anchor=north west,text=milcGradeText,opacity=.85] at ([xshift=2.0cm,yshift=-2.35cm]current page.north west)
    {\sffamily\bfseries\fontsize{8.5}{10}\selectfont SERIE GUÍAS · TECNOLOGÍA E INFORMÁTICA};
  \node[anchor=north east,text=milcGradeText,opacity=.85,align=right] at ([xshift=-2.0cm,yshift=-1.55cm]current page.north east)
    {\sffamily\bfseries\fontsize{9}{12}\selectfont I.E. Sor María Juliana\\Cartago, Valle del Cauca};

  \node[anchor=west,text=milcGradeText] (gradonum) at ([xshift=1.85cm,yshift=-7.1cm]current page.north west)
    {\sffamily\bfseries\fontsize{112}{112}\selectfont 8};
  % El circulito de grado (°) solo aplica a las guias de grado («11°»); en
  % Bebras y Semillero el numero grande es un momento/modulo, no un grado.
  % Se posiciona relativo al nodo `gradonum`, no en coordenadas absolutas.
  \draw[milcGradeText,line width=1.6pt]
    ([xshift=2.0mm,yshift=-4.5mm]gradonum.north east) circle (.15cm);

  \node[anchor=west,text=milcGradeText,text width=11.4cm,align=left] at ([xshift=1.5cm,yshift=0.5cm]gradonum.east)
    {
     {\sffamily\bfseries\fontsize{9}{11}\selectfont\MakeUppercase{Período 1 · Análisis de datos con phronesis}}\\[3mm]
     {\sffamily\bfseries\fontsize{21}{25}\selectfont\setlength{\baselineskip}{27pt}Preguntar antes\\[4pt]de calcular ---\\[4pt]phronesis con datos}\\[3.5mm]
     {\sffamily\bfseries\fontsize{15}{18}\selectfont Guía 1-8-TIC}
    };
\end{tikzpicture}

\vspace*{9.4cm}
\vfill

{\sffamily\bfseries\fontsize{8}{10}\selectfont\textcolor{milcGradeDark}{LO QUE TE LLEVAS HOY}}

\begin{tcolorbox}[enhanced,colback=white,colframe=milcNegro,arc=2mm,boxrule=1.6pt,
  left=5mm,right=5mm,top=4mm,bottom=4mm,before skip=3pt,after skip=12pt,
  fontupper=\RaggedRight\fontsize{11}{15.5}\selectfont]
Una ficha con las cinco preguntas respondidas para una hoja de datos real del colegio, la decisión de qué vas a calcular y qué dejas fuera, y un párrafo de cierre que diga cuántas filas hay, de qué periodo son y a quién no describen.
\end{tcolorbox}

\begin{tcolorbox}[enhanced,colback=milcGris,colframe=milcGris,arc=2mm,boxrule=0pt,
  left=5mm,right=5mm,top=3.5mm,bottom=3.5mm,after skip=12pt]
{\sffamily\bfseries\fontsize{8}{10}\selectfont\textcolor{milcNegro!55}{ESTA GUÍA ES DE}}\\[3mm]
\begin{tabularx}{\linewidth}{X p{3.2cm} p{3.2cm}}
\linead{\linewidth} & \linead{3.0cm} & \linead{3.0cm}\\[-1mm]
{\fontsize{8.5}{10}\selectfont\textcolor{milcNegro!55}{Nombre completo}} &
{\fontsize{8.5}{10}\selectfont\textcolor{milcNegro!55}{Grupo}} &
{\fontsize{8.5}{10}\selectfont\textcolor{milcNegro!55}{Fecha}}\\
\end{tabularx}
\end{tcolorbox}

\vfill

\noindent\textcolor{milcLinea}{\rule{\linewidth}{1pt}}\\[2mm]
{\sffamily\bfseries\fontsize{9.5}{12}\selectfont Autor: Dr. Álvaro Cárdenas Orozco}\\[1mm]
{\sffamily\fontsize{8.5}{11}\selectfont\textcolor{milcNegro!55}{Metodología MILC · apertura ancestral · pensamiento computacional · triángulo Dussel–estoicismo–Floridi}}

\newpage

\section*{Apertura: Preguntar antes de calcular --- phronesis con datos}

\begin{softbox}{milcGradeDark}{milcGradeSoft}{Saber ancestral}
En muchas tiendas de barrio de Cartago todavía existe el \textbf{cuaderno del fiado}. Cada renglón guarda cuatro datos: quién llevó, qué día, qué llevó y cuánto vale. Cuando la familia viene a pagar, la tendera no suma de una. Primero \textbf{le pregunta al cuaderno}: ¿esta letra es mía o de mi hijo?, ¿este cinco es un cinco o un tres?, ¿falta el sábado?, ¿estoy cobrando la quincena o el mes? Solo después saca la cuenta. Una socióloga que estudió el fiado en un pequeño negocio de Cali encontró que fiar es un pacto de confianza entre vecinos, sostenido por un cuaderno que se revisa y se tacha (Martínez Benavides, 2021). «Miremos el cuaderno» resuelve la discusión que la memoria sola no resuelve. Ahí está lo importante: \textbf{el registro vale porque cualquiera puede ponerlo a prueba}. Pero el cuaderno tiene una \textbf{cara de exclusión}: a quien la tendera no conoce, no le fía. Solo cuenta a los que ya estaban adentro. Preguntar antes de sumar y saber quién falta en la lista: eso es lo que hoy vas a hacer con una hoja de datos.\par\smallskip{\footnotesize\textcolor{milcNegro!70}{\textbf{Fuente.} Martínez Benavides, A. (2021). Circuitos crediticios: fiado y trabajo relacional en un pequeño negocio en Cali, Colombia. \emph{Estudios Sociológicos de El Colegio de México, 39}(116), 467--494.}}
\end{softbox}

\begin{softbox}{milcTurquesa}{milcGris}{Saber contemporáneo}
Abres la hoja de notas de 8-A en Excel: 30 filas, 6 columnas. La tentación es sacar el promedio de una. Pero tres celdas dicen «ausente», una nota está escrita con letras y nadie sabe de qué parcial es la hoja. Si calculas sin mirar, el promedio sale, se ve exacto y está mal. Un analista hace \textbf{cinco preguntas antes de la primera fórmula}: (1) \emph{¿de dónde viene la hoja?}; (2) \emph{¿quién la recogió y cuándo?}; (3) \emph{¿qué hay en cada columna: número, texto o fecha?}; (4) \emph{¿qué falta, qué se repite, qué es imposible?}; (5) \emph{¿qué quiero responder exactamente?}. A esa costumbre de preguntar antes de actuar los griegos la llamaron \textbf{phronesis}: sabiduría práctica. La tendera la tiene con su cuaderno. Hoy la practicas tú con una hoja de datos.
\end{softbox}

\begin{softbox}{milcMostaza}{milcGris}{Pregunta-puente}
\textit{La tendera le pregunta al cuaderno antes de sumar. Cuando abras una hoja con 200 filas, ¿qué le vas a preguntar antes del primer promedio? ¿Y quién no aparece en esas filas?}
\end{softbox}

\begin{softbox}{milcVerde}{milcGris}{Saber y saber hacer}
Al terminar podrás: (1) \textbf{identificar} qué columnas, filas y datos raros tiene una hoja antes de calcular; (2) \textbf{analizar} una hoja real del colegio con las cinco preguntas; (3) \textbf{evaluar} qué se puede calcular con confianza y qué se deja fuera; (4) \textbf{explicar} con tus palabras por qué preguntar antes de calcular no es perder tiempo.
\end{softbox}



\small
\begin{tabularx}{\textwidth}{>{\bfseries}p{3.0cm}Y}
\rowcolor{milcGradePrimary!12}Autor & Dr. Álvaro Cárdenas Orozco\\
DBA & Tratamiento de datos cuantitativos con criterio ético (MEN, T\&I)\\
Referentes & Phronesis aristotélica · Floridi (ética del dato) · Estoicismo (ver lo que es)\\
Producto final & Una ficha con las cinco preguntas respondidas para una hoja de datos real del colegio, la decisión de qué vas a calcular y qué dejas fuera, y un párrafo de cierre que diga cuántas filas hay, de qué periodo son y a quién no describen.\\
\end{tabularx}
\normalsize

\puente{Este es el primer periodo de 8.º y vas a aprender a trabajar con datos en Excel. Hoy no hay fórmulas: hay una costumbre. Las nueve sesiones que siguen (tipos de datos, tablas, funciones, gráficos, mini estudio) se apoyan en lo que aprendes hoy: \textbf{mirar y preguntar antes de calcular}.}

\milcestacion{milcGradeDark}{Ruta MILC: así vamos a aprender}
\begin{tabularx}{\textwidth}{>{\centering\arraybackslash}X>{\centering\arraybackslash}X>{\centering\arraybackslash}X>{\centering\arraybackslash}X}
\phasecard{1}{milcVerde}{milcGris}{Escucha\\[-1mm]\small Observo, escucho y pregunto.} &
\phasecard{2}{milcTurquesa}{milcGris}{Entender\\[-1mm]\small Sistematizo y ordeno.} &
\phasecard{3}{milcMagenta}{milcGris}{Hacer\\[-1mm]\small Construyo y aplico.} &
\phasecard{4}{milcVino}{milcGris}{Cierre\\[-1mm]\small Evalúo y reflexiono.}
\end{tabularx}


\puente{Antes de la teoría, vas a mirar una hoja de datos real durante cinco minutos sin calcular nada. Solo mirar y anotar. La primera observación es eso: observación.}

\milcestacion{milcVerde}{Estación 1 de 3 · Escucha}

\begin{softbox}{milcVerde}{milcGris}{Escena para escuchar}
\textbf{Actividad 1 · IDENTIFICA --- Mira la hoja antes de tocarla} (15 min · individual). (1) Recibe la hoja de datos que te entrega o proyecta tu docente: notas del último parcial, tiempo de pantalla del salón o gastos de la tienda escolar. (2) Mírala 5 minutos sin calcular nada: ni promedios, ni sumas. (3) Anota qué columnas tiene y cuántas filas hay. (4) Marca al menos un dato raro: una celda vacía, un texto donde va un número, un valor imposible. (5) Escribe una pregunta que te gustaría responder con esos datos. La regla de hoy: \textbf{primero observa, después pregunta, y solo al final calcula}.
\end{softbox}

{\sffamily\bfseries\fontsize{8}{10}\selectfont\textcolor{milcNegro!55}{MARCA CUANDO LO HAYAS HECHO}}
\milccheckrow{Anoté qué columnas tiene la hoja y cuántas filas hay.}
\milccheckrow{Marqué al menos un dato raro y puedo decir por qué es raro.}
\milccheckrow{Escribí una pregunta que se puede responder con esas columnas.}

\begin{infoband}{milcVerde}


\textbf{Lo que va al cuaderno (Actividad 1 · IDENTIFICA).} \textbf{Título:} «Actividad 1 --- Primera mirada». \textbf{Formato:} ficha de 4 casillas: columnas, filas, datos raros, mi pregunta. \textbf{Extensión:} media página. \textbf{Sabes que terminaste cuando} las cuatro casillas están llenas, puedes señalar con el dedo un dato raro y decir por qué lo es, y tu pregunta se responde con esas columnas y no con otras.
\end{infoband}

\puente{Ya miraste la hoja. Ahora vas a entender las \textbf{cinco preguntas} que la tendera le hace a su cuaderno antes de sumar, y a hacérselas tú a tus datos con tu pareja.}

\milcestacion{milcTurquesa}{Estación 2 de 3 · Entender}

\begin{softbox}{milcTurquesa}{milcGris}{Lo que vas a entender}
\textbf{Actividad 2 · ANALIZA --- Hazle las cinco preguntas a tu hoja} (25 min · parejas). Una hoja de datos no habla sola: contesta preguntas. Estas cinco se hacen en orden, antes de la primera fórmula. (1) Con tu pareja, toma la hoja que miraste en la Actividad 1. (2) Respondan en una línea de dónde viene y quién la recogió; si no saben, escriban «no se sabe»: también es una respuesta. (3) Recorran columna por columna y anoten el tipo de dato (número, texto, fecha) y cuántas celdas vacías o raras hay. (4) Escriban la pregunta concreta que quieren responder. (5) Marquen cada columna con C (de confianza), D (dudosa) o X (de descarte) y digan por qué en cinco palabras.
\end{softbox}

\milcpaso{1}{milcTurquesa}{\textbf{Origen y autor.} Pregunta 1: \emph{¿de dónde viene la hoja?} (una encuesta, un registro del docente, una copia de otra hoja). Pregunta 2: \emph{¿quién la recogió y cuándo?} Una hoja de hace tres años responde preguntas de hace tres años. Si no sabes el origen, escríbelo: «origen desconocido» ya es un dato sobre la hoja.}
\milcpaso{2}{milcTurquesa}{\textbf{Tipo de dato.} Pregunta 3: \emph{¿qué hay en cada columna?} Números, texto o fechas. Excel suma números, no palabras: «dos mil» escrito con letras no entra en la suma y nadie te avisa. Una fecha escrita como texto tampoco se ordena. Mira columna por columna, no la hoja entera.}
\milcpaso{3}{milcTurquesa}{\textbf{Errores y vacíos.} Pregunta 4: \emph{¿qué falta, qué se repite, qué es imposible?} Una celda vacía no es un cero: si la tratas como cero, el promedio baja. Una nota de 7 en una escala de 1 a 5 es imposible. Un estudiante repetido cuenta dos veces. Anota cuántas celdas raras hay en cada columna.}
\milcpaso{4}{milcTurquesa}{\textbf{Pregunta y decisión.} Pregunta 5: \emph{¿qué quiero responder exactamente?} «¿Cómo le fue al grupo?» es vaga; «¿cuántos estudiantes sacaron menos de 3,0 en el parcial de marzo?» se puede responder. Con la pregunta clara, marca cada columna: \textbf{C} de confianza (origen claro, tipo correcto, pocos vacíos), \textbf{D} dudosa (algo no cuadra) o \textbf{X} de descarte (mejor no usarla).}

\begin{drawbox}{milcTurquesa}{Las cinco preguntas, en una mirada}
\textbf{1 Origen:} ¿de dónde viene? \\[1mm] \textbf{2 Autor y fecha:} ¿quién y cuándo? \\[1mm] \textbf{3 Tipo de dato:} ¿número, texto o fecha? \\[1mm] \textbf{4 Errores y vacíos:} ¿qué falta o es imposible? \\[1mm] \textbf{5 Pregunta:} ¿qué quiero responder?
\end{drawbox}

\begin{softbox}{milcMostaza}{milcGris}{Errores comunes y su corrección}
(1) \textbf{Saltar al cálculo}: abrir Excel y sacar promedios sin haber mirado la hoja. (2) \textbf{Creer que una celda vacía es un cero}: baja el promedio en silencio. (3) \textbf{Sumar texto disfrazado de número}: «dos mil», «3,5 aprox.», «N/A». (4) \textbf{Pensar que más filas es más confiable}: 200 filas del mismo salón siguen hablando solo de ese salón. (5) \textbf{No decir qué dejaste fuera}: presentar el resultado como si la hoja fuera perfecta.
\end{softbox}

\begin{infoband}{milcTurquesa}
\guiaDiagrama{assets/1-8/mirar-preguntar-decidir.tex}{El orden importa: mirar, preguntar y decidir van antes de cualquier fórmula. La tendera hace lo mismo con su cuaderno antes de sumar.}

\textbf{Lo que va al cuaderno (Actividad 2 · ANALIZA).} \textbf{Título:} «Actividad 2 --- Las cinco preguntas». \textbf{Formato:} tabla de 5 filas (una por pregunta) y 2 columnas (pregunta / respuesta), más la lista de columnas con su letra C, D o X. \textbf{Extensión:} media página; cada respuesta de 1 a 2 renglones. \textbf{Sabes que terminaste cuando} las cinco preguntas tienen respuesta (aunque alguna sea «no se sabe»), cada columna tiene su letra con una razón, y tu pareja puede decir, sin preguntarte nada, si confiaría en un promedio de esa hoja.
\end{infoband}

\puente{Ya sabes qué tan confiable es tu hoja. Toca decidir: qué vas a calcular, qué dejas fuera y por qué. Lo escribes tú, con honestidad.}

\milcestacion{milcMagenta}{Estación 3 de 3 · Hacer}

\begin{softbox}{milcMagenta}{milcGris}{Cómo vas a construir tu producto}
\textbf{Actividad 3 · EVALÚA --- Decide qué calcular y qué dejar fuera} (25 min · individual). Ya sabes qué tan confiable es tu hoja. Ahora decides, y lo escribes tú. (1) Relee las respuestas de las cinco preguntas; tu pareja puede opinar, pero escribes tú. (2) Elige la pregunta que sí se puede responder con confianza y subráyala. (3) Escribe qué vas a calcular (promedio, suma, conteo, mínimo, máximo) y por qué sirve para tu pregunta. (4) Escribe qué dejas fuera y por qué: una columna dudosa, celdas vacías, un dato que no entiendes. (5) Cierra con un párrafo de 4 renglones: cuántas filas hay, de qué periodo son y a quién sí y a quién no describen.
\end{softbox}

\milcpaso{1}{milcMagenta}{\textbf{Tu ficha tiene 4 piezas.} (1) La pregunta elegida, subrayada. (2) «Voy a calcular»: uno o dos cálculos, cada uno con su razón. (3) «Dejo fuera»: al menos una cosa, con su razón. (4) Un párrafo de cierre que diga cuántas filas hay, de qué periodo son y a quién no describen.}
\milcpaso{2}{milcMagenta}{\textbf{Calcular, calcular con aviso o no calcular.} Calcula cuando la columna es C: origen claro, tipo correcto, pocos vacíos. Calcula con aviso cuando es D: sirve, pero escribes el límite al lado («solo 12 estudiantes», «solo la semana pasada»). No calcules cuando es X: muchos vacíos, tipo equivocado, origen desconocido.}
\milcpaso{3}{milcMagenta}{\textbf{Honestidad antes que apariencia.} Una ficha que dice «esta hoja tiene 30 filas de un solo salón y no describe al colegio» es mejor trabajo que una que finge certeza. Los datos siempre tienen huecos; el oficio consiste en mostrarlos, no en taparlos. Quien tapa los huecos engaña; quien los muestra construye confianza.}
\milcpaso{4}{milcMagenta}{\textbf{Antes de entregar, léelo en voz alta.} Si suena defensivo («los datos estaban mal, no es culpa mía»), reescribe: describe la hoja, no te disculpes. Si suena seguro de más («esto prueba que…»), agrega el límite. El tono correcto es el de la tendera: «miremos el cuaderno».}

\begin{drawbox}{milcMagenta}{Antes de entregar}
\checkbox Las cinco preguntas tienen respuesta de 1 a 2 renglones. \\[1mm] \checkbox La pregunta elegida está subrayada y tiene una razón. \\[1mm] \checkbox Cada cálculo tiene su razón escrita al lado. \\[1mm] \checkbox Hay al menos una cosa que dejo fuera, con su razón. \\[1mm] \checkbox El párrafo de cierre dice cuántas filas hay y de qué periodo son. \\[1mm] \checkbox Lo leí en voz alta y suena honesto, no defensivo.
\end{drawbox}

\begin{softbox}{milcMagenta}{milcGris}{Plantilla de trabajo}
\textbf{Pregunta elegida.} \textit{La que sí se puede responder.} \\[1mm] \textbf{Voy a calcular.} \textit{Uno o dos cálculos, cada uno con su razón.} \\[1mm] \textbf{Dejo fuera.} \textit{Qué no uso y por qué.} \\[1mm] \textbf{Cierre.} \textit{Cuántas filas, de qué periodo, a quién no describen.}
\end{softbox}

\begin{infoband}{milcMagenta}


\textbf{Lo que va al cuaderno (Actividad 3 · EVALÚA).} \textbf{Título:} «Actividad 3 --- Decido qué calcular». \textbf{Formato:} ficha de 4 piezas: pregunta elegida, calculo, dejo fuera, párrafo de cierre. \textbf{Extensión:} 1 página. \textbf{Sabes que terminaste cuando} cada cálculo tiene una razón al lado, hay al menos una cosa que dejas fuera con su porqué, el párrafo dice cuántas filas hay y de qué periodo son, y lo leíste en voz alta sin sonar defensivo.
\end{infoband}

\puente{Lo que entregas hoy: la ficha de las cinco preguntas, la decisión de qué calcular y qué no, y un párrafo de cierre con los límites de tu hoja. La prueba de calidad: un compañero lee tu ficha y sabe si confiar o no en tu cálculo.}

\milcestacion{milcMostaza}{Producto de la sesión}
\textbf{Ficha de las cinco preguntas + decisión de qué calcular + cierre honesto}.
\begin{softbox}{milcMostaza}{milcGris}{Criterios de calidad}
(1) Las cinco preguntas tienen respuesta de 1 a 2 renglones cada una. (2) La pregunta elegida está subrayada y tiene una razón escrita. (3) Cada cálculo propuesto tiene su razón al lado. (4) Hay al menos una cosa que se deja fuera, con su porqué. (5) El párrafo de cierre dice cuántas filas hay y de qué periodo son. (6) La ficha cabe en una página y se leyó en voz alta.
\end{softbox}

\puente{Antes de cerrar, mira lo que hiciste desde cinco lados. Preguntar antes de calcular es respetar los datos y a las personas que hay detrás; saltar al cálculo es la forma más rápida de engañarse.}

\milcestacion{milcVino}{Evaluación liberadora · cinco dimensiones}

\begin{softbox}{milcVino}{milcGris}{Sentido de la evaluación}
Evaluar no es solo revisar una nota. Es reconocer qué aprendí, qué cambió en mí, qué emociones manejé, cómo me relaciono con la ciudad y la comunidad, y qué le dejaré a quien viene detrás.
\end{softbox}

\textbf{Mi reflexión liberadora en cinco dimensiones.} Completa cada frase iniciada:

\small
\begin{tabularx}{\textwidth}{>{\bfseries\RaggedRight\arraybackslash}p{3.4cm}Y>{\RaggedRight\arraybackslash}p{4.6cm}}
\rowcolor{milcVino}\color{white}Dimensión & \color{white}\bfseries Mi respuesta (frame starter) & \color{white}\bfseries Algo que descubrí\\
1. Personal & Lo que más cambió en mí fue \linead{6.5cm} & \linead{4.2cm}\\
2. Emocional & La emoción más fuerte fue \linead{2.5cm} y la regulé \linead{3cm} & \linead{4.2cm}\\
3. Ciudadana & En democracia esto importa porque \linead{6cm} & \linead{4.2cm}\\
4. Local (Cartago) & En mi barrio sería útil para \linead{6.5cm} & \linead{4.2cm}\\
5. Intergeneracional & A mi abuela/abuelo le diría \linead{6.5cm} & \linead{4.2cm}\\
\end{tabularx}
\normalsize

\begin{infoband}{milcVino}
\textbf{Para la próxima sustentación.} Vuelve a esta tabla en seis meses. Verás que las respuestas cambian — no porque lo aprendido cambie, sino porque tú cambias. La evaluación liberadora es una conversación contigo a través del tiempo.
\end{infoband}


\puente{Para terminar, tres voces sobre mirar bien. Dussel, en \emph{Filosofía de la liberación}, sobre quién queda fuera del sistema; Marco Aurelio, en las \emph{Meditaciones}, sobre cambiar de opinión cuando te muestran el error; Floridi, en su texto sobre big data, sobre cuándo un patrón pequeño significa algo.}

\milcestacion{milcGradeDark}{Triángulo de pensamiento · cierre filosófico}

\begin{softbox}{milcVino}{milcGris}{Dussel · lente del nosotros}
\textit{«El otro se revela realmente como otro\ldots como el pobre, el oprimido; el que, a la vera del camino, fuera del sistema, muestra su rostro sufriente y sin embargo desafiante.»}

{\footnotesize\itshape\color{milcNegro!70}--- Enrique Dussel · Filosofía de la liberación (1977)}

La hoja solo cuenta a quien alguien anotó. El que no vino el día de la encuesta, el que no tiene celular para reportar su tiempo de pantalla, el que la tendera no conoce: ninguno está en las filas. Antes de calcular, pregunta quién quedó \textbf{a la vera del camino}. Un promedio que no sabe a quién dejó fuera repite la exclusión sin darse cuenta.

\textbf{¿Quién no aparece en mi hoja? ¿Qué cambiaría en mi cálculo si estuviera?}
\end{softbox}
\linead{\linewidth}\\[1mm]
\linead{\linewidth}

\begin{softbox}{milcMostaza}{milcGris}{Estoicismo · Marco Aurelio · Meditaciones VI, 21 (c. 175 d.C.) · cuidado interior}
\textit{«Si alguien puede refutarme y probar de modo concluyente que pienso o actúo incorrectamente, de buen grado cambiaré de proceder. Pues persigo la verdad, que no dañó nunca a nadie; en cambio, sí se daña el que persiste en su propio engaño e ignorancia.»}

{\footnotesize\itshape\color{milcNegro!70}--- Marco Aurelio · Meditaciones VI, 21 (c. 175 d.C.)}

Si tu pareja te muestra que una columna está mal, cambias la decisión y ya: no es una derrota, es lo que hace el que persigue la verdad. Persistir en un promedio que sabes flojo, porque ya lo escribiste o porque te da la razón, es el engaño del que habla el emperador. \textbf{Ver los datos como son} vale más que tener la razón.

\textbf{¿Estoy leyendo los datos como son, o como me gustaría que fueran para darme la razón?}
\end{softbox}
\linead{\linewidth}\\[1mm]
\linead{\linewidth}

\begin{softbox}{milcTurquesa}{milcGris}{Floridi · lente de la infoesfera}
\textit{«Los pequeños patrones solo pueden ser significativos si se agregan correctamente, se comparan y se procesan a tiempo.»}

{\footnotesize\itshape\color{milcNegro!70}--- Luciano Floridi · Big data and their epistemological challenge (2012)}

Un promedio es un patrón pequeño. Solo vale si sumaste bien (sin textos disfrazados de números ni vacíos contados como ceros), si lo comparaste con algo (otro salón, otro mes) y si los datos no están vencidos. Las cinco preguntas son la forma de \textbf{agregar correctamente} antes de creerle a un número.

\textbf{¿Mi cálculo agrega bien, compara con algo y llega a tiempo? ¿O es un número solo?}
\end{softbox}
\linead{\linewidth}\\[1mm]
\linead{\linewidth}

\begin{infoband}{milcNegro}
\textbf{Nota para el docente.} Las tres son citas textuales y llevan su referencia debajo. Puedes remitir a la obra en clase.
\end{infoband}

\milcestacion{milcVino}{Mi autoevaluación · marca una casilla por fila}

{\small
\setlength{\extrarowheight}{2.2pt}
\begin{tabularx}{\textwidth}{>{\RaggedRight\arraybackslash\bfseries}p{3.0cm}YYY}
\rowcolor{milcVino}\color{white}\bfseries Criterio & \color{white}\bfseries Lo logré & \color{white}\bfseries En proceso & \color{white}\bfseries Necesito apoyo\\[.6mm]
Comprensión del tema & \milcopcion{Explico las ideas con mis palabras.} & \milcopcion{Reconozco las ideas, falta precisión.} & \milcopcion{Necesito repasar lo básico.}\\[.8mm]
\rowcolor{milcGris}Pensamiento computacional & \milcopcion{Usé los pilares al armar mi producto.} & \milcopcion{Usé algunos pilares.} & \milcopcion{Necesito releer la estación 2.}\\[.8mm]
Aplicación y producto & \milcopcion{Mi producto cumple los criterios.} & \milcopcion{Está armado, con ajustes pendientes.} & \milcopcion{Aún está en construcción.}\\[.8mm]
\rowcolor{milcGris}Reflexión 5 dimensiones & \milcopcion{Respondí cada una con honestidad.} & \milcopcion{Respondí algunas con detalle.} & \milcopcion{Mis respuestas son muy generales.}\\[.8mm]
Triángulo de pensamiento & \milcopcion{Conecté las 3 voces con mi trabajo.} & \milcopcion{Conecté al menos una.} & \milcopcion{No las relaciono con mi proceso.}\\[.8mm]
\end{tabularx}}

\vspace{3mm}
\textbf{Hoy entendí que:}\\[1mm]
\linead{\linewidth}\\[1mm]
\linead{\linewidth}

\textbf{Mi compromiso para la próxima sesión es:}\\[1mm]
\linead{\linewidth}\\[1mm]
\linead{\linewidth}

\begin{softbox}{milcMostaza}{milcGris}{Cita de despedida · Marco Aurelio}
\textit{``El obstáculo en el camino se vuelve el camino. Lo que estorba al camino se vuelve camino.''}\\[1mm]
Cada error de hoy, bien observado, se convierte en pieza del camino que pisarás mañana.
\end{softbox}

\small
\begin{tabularx}{\textwidth}{>{\bfseries}p{3.6cm}Y}
\rowcolor{milcGradeDark!12}Nombre completo: & \linead{10cm}\\
Fecha: & \linead{4cm} \quad Curso/grupo: \linead{4cm}\\
Mi firma: & \linead{9cm}\\
\end{tabularx}
\normalsize

\begin{infoband}{milcGradeDark}
\textbf{Para la próxima sesión.} Llega con tu hoja elegida y las cinco preguntas respondidas. En la sesión 2 vas a clasificar esos datos por tipo en Excel: número, texto, fecha, moneda. Lo de hoy es la base; lo de la próxima semana se monta encima.
\end{infoband}

\end{document}
